\usepackage[acronym]{glossaries}


\makeglossaries

% -------------------------
% Acronyms
% -------------------------

\newacronym{abs}{\(B\)}{Automation Bias Score}
\newacronym{mean-abs}{\(\bar{B}\)}{Mean Automation Bias Score}
\newacronym{h1}{\(H_1\)}{Hypothesis 1: Short Term Effect}
\newacronym{h2}{\(H_2\)}{Hypothesis 2: Long Term Effect}
\newacronym{G}{\(G\)}{possible participant groups}
\newacronym{C}{\(C\)}{possible case vignettes}
\newacronym{T}{\(T\)}{possible truth values}
\newacronym{R}{\(R\)}{possible rounds for participants}
\newacronym{S}{\(S\)}{possible sessions for participants}
\newacronym{S2}{\(S^2\)}{possible session combinations for participants}
\newacronym{P}{\(P\)}{possible participants}

\newacronym{ai}{AI}{Artificial Intelligence}
\newacronym{ueq-plus}{UEQ+}{User Experience Questionnaire Plus}
\newacronym{thi}{THI}{Technische Hochschule Ingolstadt}
\newacronym[
    description={A web application that uses modern web technologies to provide an app-like experience, including offline functionality, responsiveness, and installability across different devices and platforms}
]{pwa}{PWA}{Progressive Web App}

% -------------------------
% Glossary Entries
% -------------------------


\newglossaryentry{pseudo-random-seed}
{
name=pseudo-random seed,
plural=pseudo-random seeds,
description={A fixed numerical value used to ensure reproducibility of pseudo-random participant allocation}
}

